\section{Preliminaries} \label{sec:prelim}



\subsection{The Set-Intersection Problem}

\begin{definition}
    The \textbf{$m$-bit Set-Intersection} problem is a communication problem
    for two players Alice and Bob.
    Alice is given a subset $X \subseteq [m]$ and Bob is given a subset
    $Y \subseteq [m]$, with the promise that $|X \cap Y| = 1$.
    Their goal is to learn the point of intersection $\{t\} = X \cap Y$. 
\end{definition}

Let us define a distribution $\cD_\SI$ over inputs for Set-Intersection. This is the positive case of Razborov's original hard distribution for disjointness \cite{Razborov92}.

\begin{tbox}
    \textbf{Distribution $\cD_\SI$}~\cite{AssadiR20} on sets
    $(X, Y)$ from the universe $\cX=[m]$: 

    \begin{itemize}
    \item Uniformly sample two disjoint sets $X', Y' \subseteq [m]$ of size
      $\frac m 4 - 1$.
    \item Sample $\estar \in [m]\setminus(X' \cup Y')$ uniformly at random
      and define $X = X' \cup \Set{\estar}$ and $Y = Y' \cup \Set{\estar}$.
\end{itemize}
\end{tbox}

\begin{proposition}\label{prop:SI-hard}
  The internal information complexity of \emph{finding} the intersecting element
  in a $\setint_m$ instance sampled from $\dsi$ is $\Omega(m)$.
\end{proposition}

\Cref{prop:SI-hard} is not directly implied by standard set-intersection lower
bounds, since they use distributions where $X$ and $Y$ never intersect.
To obtain this proposition, one can (for example) apply the result
of~\cite{JayramKS03} on the information cost of set disjointness protocols on
intersecting distributions (see~\cite{AssadiCK19} for more details).

\subsection{The (Boolean) Hidden Pointer Chasing Problem}

The Hidden Pointer Chasing ($\hpc$) problem is a composition of Pointer Chasing and Set-Intersection, defined as follows.

\begin{definition}\label{prob:hpc}
The \textbf{Hidden Pointer Chasing ($\HPC$)} problem is a communication problem
with four players $P_A, P_B, P_C$, and $P_D$. 
Let $\cX := \{x_1,...,x_m\}$ and $\cY := \{y_1,...,y_m\}$ be two disjoint universes. 
    \begin{itemize}
        \item For any $x\in \cX$, $P_A$ and $P_B$ are given an instance $(A_x,B_x)$ of \setint over the universe $\cY$ where $A_x \cap B_x = \{t_x\}$ for a single target element $t_x \in \cY$. We define $A := \{A_{x_1} ,\ldots, A_{x_m} \}$ and $B:=\{B_{x_1} ,\ldots, B_{x_m} \}$ as the whole input to $P_A$ and $P_B$, respectively.
        
        \item For any $y\in \cY$, $P_C$ and $P_D$ are given an instance $(C_y,D_y)$ of $\setint$ over the universe $\cX$ where $C_y \cap D_y = \{t_y\}$ for a single target element $t_y \in \cX$. We define $C := \{C_{y_1} ,\ldots, C_{y_m} \}$ and $D:=\{D_{y_1} ,\ldots,D_{y_m} \}$ as the whole input to $P_C$ and $P_D$, respectively.
        
        \item We define two mappings $f_{AB} : \cX \to \cY$ and $f_{CD} : \cY \to \cX$ such that:
        \[
            (a) ~ \textnormal{for any $x\in \cX$, $f_{AB}(x) = t_x$;} \qquad \qquad (b) ~ \textnormal{for any $y\in \cY$, $f_{CD}(y) = t_y$}.
         \]
        \item Let $x_1 \in \cX$ be an arbitrary fixed element known to all players. The pointers $z_0,z_1,z_2,\ldots$ are defined inductively as 
        $$z_0 := x_1, ~~z_1 := f_{AB}(z_0),~~ z_2 := f_{CD}(z_1), ~~z_3 := f_{AB}(z_2),~\ldots.$$
    \end{itemize}
For any integers $m,r \geq 1$, the \textnormal{\textbf{$m$-bit, $r$-step Hidden Pointer Chasing}} problem, denoted by $\hpc_{m,r}$ is defined as the communication problem of finding the pointer $z_r$. Meaning, we wish to output the index $i \in [m]$ such that $z_r = x_i$ (when $r$ is even) or $z_r = y_i$ (when $r$ is odd). 

We also define a Boolean variant of $\hpc$ as follows. In the \textnormal{\textbf{$m$-bit, $r$-step \underline{Boolean} Hidden Pointer Chasing}} problem, denoted by $\bhpc_{m,r}$, the goal is to output the bit $b(z_r):=i \mod 2$, where $z_r = x_i$ or $y_i$ depending on if $r$ is even or odd respectively. 
\end{definition}

\subsection{The (Boolean) Multilayer Hidden Pointer Chasing Problem}

The Multilayer Hidden Pointer Chasing problem ($\mhpc$) combines
$r$ instances of $\hpc$ as its input in the following manner.


\begin{definition}
  The \textbf{Multilayer Hidden Pointer Chasing} ($\mhpc$) problem is a
  communication problem with four players $P_A, P_B, P_C$ and $P_D$. Let $m,r$ be integers $\geq 1$. The $m$-bit, $r$-layer \mhpc problem, denoted by $\mhpc_{m,r}$ is defined as follows.
  \begin{itemize}
    \item For each $i \in \bracket{r}$, the players receive an $\hpc$ instance
      $(A^i, B^i, C^i, D^i)$, each over the disjoint universes
      $\cX$ and $\cY$ of size $m$.
      Let $t^i_x$ denote the unique element in the intersection
      $A^i_x \cap B^i_x$, and $t^i_y$ the unique element in $C^i_y \cap D^i_y$.
    \item For each $i \in \brac{r}$, we have the maps $f^i_{AB} : \cX \to \cY$
      and $f^i_{CD} : \cY \to \cX$, defined as:
      \[
        (a) ~ \textnormal{for any $x\in \cX$, $f^i_{AB}(x) = t^i_x$;}
        \qquad \qquad
        (b) ~ \textnormal{for any $y\in \cY$, $f^i_{CD}(y) = t^i_y$}.
      \]
    \item
      The pointers $z_i$ are defined inductively as:
      \[
        z_0 := x_1, ~~z_1 := f^1_{AB}(z_0),~~ z_2 := f^2_{CD}(z_1),
        ~~z_3 := f^3_{AB}(z_2),~\ldots.
      \]
  \end{itemize}
  % As before, the goal for the players is to compute $z_r$.

  The goal of the problem is to find $z_r$, i.e., the index $i \in [m]$ such that $z_r = x_i$ (when $r$ is even) or $z_r = y_i$ (when $r$ is odd). 

Analogous to \hpc, we define the Boolean variant of $\mhpc$ as follows. In the \textnormal{\textbf{$m$-bit, $r$-layer \underline{Boolean} \mhpc}} problem, denoted by $\bmhpc_{m,r}$, the goal is to output the bit $b(z_r):=i \mod 2$, where $z_r = x_i$ or $y_i$ depending on if $r$ is even or odd respectively. 
\end{definition}

We remark that $C^j, D^j$ for odd $j$, and $A^j, B^j$ for even $j$, are redundant (they do not affect the output of the function). They are kept for notational convenience and ease of presentation. 

\subsection{Misaligned \texorpdfstring{$r$-Round}{r-Round} Protocols}

We are interested in bounded-round protocols for $\hpc$, $\bhpc$, $\mhpc$ and $\bmhpc$. In each round, two of the players speak with each other ($P_A$ with $P_B$, or $P_C$ with $P_D$), but not with the other pair of players, with a message between pairs of players triggering the next round.

\begin{definition}[Round of $\hpc_{m,r}$]\label{def:round}
An $r$-round protocol is a protocol where the communication is organized in rounds $1, 2, \ldots, r$. A given pair of players --- $P_A$ and $P_B$ or $P_C$ and $P_D$ --- communicate with each other in the odd rounds (the first, third, etc rounds) and the other pair of players communicate with each other in the even rounds. Each round ends with the speaking pair of players sending a message to the other pair of players, and this marks the beginning of the next round.
\end{definition}

It is easy to see that one can solve any of the above HPC variants by $r$-round protocols communicating $O(r m)$ bits, if $P_A$ and $P_B$ speak in the odd rounds, and $P_C$ and $P_D$ speak in the even rounds. We will show lower-bounds for \textit{misaligned} $r$-round protocols where the \textit{wrong} pair of players speak:

\begin{definition}
We call an $r$-round protocol \textit{misaligned} if the pair of players $P_C$ and $P_D$ speak in the odd rounds, and $P_A$ and $P_B$ speak in the even rounds.
\end{definition}

The above definitions naturally suggest that, in an $r$-round protocol, we should think of $z_j$ as the pointer that the players wish to know by the end of round $j$. In a misaligned protocol, however, the players speaking in round $j$ do not hold the Set-Intersection instance that needs to be solved to find $z_j$.

\subsection{Notation and Terminology}

\mypar{Notation for Protocols} We use the notation $\prot_\SI$, $\prot_\bhpc$,
and $\prot_\bmhpc$, to denote protocols for Set-Intersection, $\bhpc$ and
$\bmhpc$, respectively. 

In general, we will use the upright $\rA$ to denote a random variable
associated with $A$.

Now consider a situation where we choose inputs $(A, B, C, D)$ from some
distribution on inputs to $\hpc_{m,r}$.
We then let $\rZ_1, \ldots , \rZ_r$ denote random variables with values
$z_1, \ldots, z_r$ of \Cref{prob:hpc}.
We use the variable $j \in [r]$ to index the rounds of a
$r$-round protocol.
For any $j \in [r]$, we define $\Prot_j$ as the set of all messages
communicated in round $j$, and more generally $\Prot$ denotes the entire
protocol transcript. 

For any $x \in \cX$ and $y \in \cY$, we define random variables
$\rT_x \in \cY$ and $\rT_y \in \cX$ giving the intersecting elements of the
Set-Intersection instance $(\rA_x, \rB_x)$ and $(\rC_y, \rD_y)$, respectively.
We define $\rE_j = (\rPi_1, \ldots, \rPi_{j-1}, \rZ_1, \ldots, \rZ_{j-1})$ to
include all messages exchanged before round $j$, together with the entire path
followed by the pointers before round $j$.
Note that during round $j$, the players wish to compute the intersection
$A_{z_j} \cap B_{z_j}$ (or $C_{z_j} \cap D_{z_j}$), so knowing $\rE_j$ should
not help them to do that. (Note that $\rE_1$ is empty.)

On an input $(A, B, C, D)$ to $\bmhpc_{m,r}$, we define the notation
$Z_i$, $j$, $\Prot_j$, $T^j_x$, $T^j_y$ analogously.

If we have random variables $\rA, \rB$, we will broadly use the notation
$\mu(\rA \mid \rB)$ to mean ``the distribution of $\rA$, conditioned on the
value of $\rB$ that occurred''.
So, for example, if $\rB$ is constant we may write $\mu(\rA)$, and this
is some fixed distribution.
But if $\rB$ is non-constant, the distribution $\mu(\rA \mid \rB)$ may vary
depending on $\rB$, and so $\mu(\rA \mid \rB)$ is itself also a random
variable, which depends on $\rB$ and can be determined by knowing $\rB$ alone
(even without knowing $\rA$).

\mypar{Other notation}
An input graph to a streaming algorithm will usually be called $G=(V,E)$ and $n$ will denote $|V|$. The degeneracy of graph $G$ will be denoted by $\kappa(G)$, or simply $\kappa$ if the graph is clear from the context. The notation ``$\log x$'' stands for $\log_2 x$. For a positive integer $a$, the notation $[a]$ denotes the set $\{1,2,\ldots,a\}$ and for integers $a<b$, $[a,b]$ denotes the set $\{a,a+1,\ldots,b\}$.
 


\subsection{Positive Triangular Discrimination} \label{sec:td}

\begin{definition}[Positive triangular discrimination]
    Given two distributions $\mu$ and $\nu$ on a sample space $\Omega$, the \textit{positive triangular discrimination} 
    % $\Delta(\mu, \nu)$ is defined as the following measure of distance between $\mu$ and $\nu$:
    % \[
    % \Delta(\mu, \nu) = \sum_{x \in x}\frac{(\mu(x) - \nu(x))^2}{\mu(x) + \nu(x)}.
    % \]
    % The positive triangular discrimination 
    $\Lambda(\mu,\nu)$ is:
    \[
    \Lambda(\mu, \nu) = \sum_{\substack{x \in \Omega\\\mu(x) > \nu(x)}}\frac{(\mu(x) - \nu(x))^2}{\mu(x) + \nu(x)}.  %\leq \Delta(\mu, \nu).
    \]
\end{definition}

In this paper, we work exclusively with \textit{positive} triangular discrimination. Below we state some of the useful properties of positive triangular discrimination. We defer the proofs of the following propositions to \Cref{sec:td-props}.


\begin{restatable}{proposition}{lambdatvd}\label{prop:prop-lambda}
$0 \le \frac{\|\mu - \nu\|^2_1} 8\leq \Lambda(\mu, \nu)\leq \frac{\|\mu - \nu\|_1} 2 \le 1$. 
\end{restatable}


The following property is already implicit in \cite{Yehudayoff20}, but making it explicit now simplifies the presentation later. Note that an analogous equation for total variation distance, i.e., $\E_\mu[f(x)] \le \|\mu-\nu\|_1\cdot \max_{x \in x} f(x) + \cdot \E_{\nu}[f(x)]$ can be obtained almost from the definition. The following statement, however, needs a bit more work.

\begin{restatable}{proposition}{lambdaloss}\label{itm:lambda-loss}
For any function $f:\Omega\to\bbR_{\ge 0}$,
\[ \E_\mu[f(x)] \le \Lambda(\mu,\nu)\cdot \max_{x \in \Omega} f(x) + 6 \cdot \E_{\nu}[f(x)]. \]
\end{restatable}



\noindent
We will also make use of the convexity of $\Lambda$.

\begin{restatable}{proposition}{convex}\label{prop:convex}
Positive triangular discrimination $\Lambda(\mu,\nu)$ is convex as a joint function of the two distributions $\mu$ and $\nu$. 
\end{restatable}



\subsection{Degeneracy and \texorpdfstring{$k$-Cores}{k-Cores}}

The definition of a $k$-core and the degeneracy $\kappa(G)$ of a graph $G$ can be found in \Cref{sec:degen-intro}. Here, we define some related terms that come up often in the paper. We also state some basic facts about them that we later use in our proofs. 


Given an ordering $\sigma$ of vertices, we say $v\prec_\sigma u$ if $u$ appears (somewhere) after $v$ in $\sigma$. 

\begin{definition}[Outdegree with respect to an ordering]\label{def:odegordering}
    For a graph $G=(V,E)$. given an ordering $\sigma$ of vertices in $V$, the outdegree of a vertex $v$ with respect to $\sigma$, denoted by $\odeg_{\sigma}(v)$, is the number of neighbors of $v$ that appear after it in $\sigma$. Formally,
$$\odeg_{\sigma}(v) = |\{u\in N(v): v \prec_\sigma u\}|$$ 
\end{definition}

We can also think of it as orienting all edges $(x,y)$ from $x$ to $y$ if $x\prec_\sigma y$. Then $\odeg_\sigma(v)$ is the outdegree of $v$ in the resultant oriented graph. 


\begin{definition}[$k$-ordering]\label{def:k-ordering} 
    Given a graph $G=(V,E)$, an ordering $\sigma$ of vertices in $V$ is called a $k$-ordering if $\forall v\in V: \odeg_\sigma(v)\leq k$.   
\end{definition}

The following fact connects vertex orderings to degeneracy. 

\begin{fact}[\cite{MatulaB83}]\label{fact:min-k-ord}
    The minimum number $k$ for which a graph $G$ admits a $k$-ordering is $\kappa(G)$.
\end{fact}


The above fact implies that to prove that a graph has degeneracy at most $k$, it suffices to exhibit a $k$-ordering of the nodes. We now formally define \emph{degeneracy ordering} of a graph.  

\begin{definition}[Degeneracy ordering] \label{def:degen-ord}
    Given a graph $G=(V,E)$, a $\kappa(G)$-ordering of $V$ is called a degeneracy ordering of $G$.
\end{definition}

A degeneracy ordering of a graph can be obtained by running the following ``peeling algorithm'' on it: starting with an empty sequence $\sigma$, recursively remove the minimum degree vertex from $G$ and append it to $\sigma$ until the graph becomes empty. Again, the maximum degree that a node has when it is peeled by this algorithm is exactly $\kappa(G)$. If not, then the peeling algorithm gives a $k$-ordering for some $k<\kappa(G)$ which violates \Cref{fact:min-k-ord}. Hence, we get the following fact. 


 
  

\begin{fact}\label{fact:peelingmin-degen}
Suppose the peeling algorithm on a graph $G$ removes its nodes in the order $v_1,\ldots,v_n$. Define $G_0:=G$ and $G_i:= G\setminus \{v_1,\ldots,v_i\}$ for $i\in [n]$. Then, $$\kappa(G)=\max_{i\in [n]} \deg_{G_{i-1}}(v_i)$$. 
\end{fact}



